\FloatBarrier
\newpage
\section{综合习题}
\label{sec:calculus-exercises}

以下题目覆盖全章，可根据问题选择已有的方法。星号题使用相同的函数与累积对象，进一步研究局部模型、统一误差及函数运算。

\begin{exercise}
\mbox{}\par\nobreak

（1）设\(f\)在\(a\)处可导。求\(\displaystyle\lim_{h\to0}\frac{f(a+3h)-f(a-2h)}h\)，并用微分解释各增量的作用。

（2）设\(f(0)=0\)，\(x\ne0\)时\(f(x)=x^2\sin(1/x)\)。判断它的连续性、可导性及导函数在零点的连续性。再与\(g(0)=0\)、\(g(x)=x\sin(1/x)\)比较，说明函数增量与差商分别需要怎样的控制。

（3）设\(f(x)=x^3+x\)，反函数为\(g\)。求\(g'(2)\)，再求\(\int_0^2g'(x)\,dx\)。第二问是否必须先求出\(g\)的显式表达式？

（4）讨论\(f_a(x)=e^x-ax\)在\(\mathbb R\)上的单调性、最小值及零点个数，并确定\(f_a(x)\ge0\)恒成立的参数范围。

（5）设\(f\)在\([0,1]\)上连续、内部可导，且\(f(0)=f(1)=0\)、\(|f'|\le2\)。估计\(|f(x)|\)以及\(\int_0^1|f(x)|\,dx\)，并构造达到积分上界的连续分段线性函数。比较这个函数与原题可导要求的关系。

（6）以2024年新课标I卷第10题的函数\(f(x)=(x-1)^2(x-4)\)为研究对象：求其单调区间与极值，讨论方程\(f(x)=k\)的不同实根个数，再求图像在\([1,4]\)上与横轴围成的面积。

（7）求\(y=e^x\)与\(y=-e^{-x}\)的公共切线。分别设两个切点，比较公共切线与公共点处相切的条件。

（8）设\(A(x)=\int_0^x(t^2-2t)\,dt\)，\(0\le x\le3\)。求其单调区间、最值与零点，并解释\(A\)的增减为何不由\(t^2-2t\)本身的单调性直接决定。

（9）设\(f\)在\([0,1]\)上连续，且\(f(x)+f(1-x)=1\)。求\(\int_0^1f(x)\,dx\)。\(\int_0^1xf(x)\,dx\)能否仅由这些条件唯一确定？举例或证明。

（10）计算\(\int_0^1 x/(1+x^2)\,dx\)、\(\int_0^{\pi/2}x\cos x\,dx\)和\(\int_{-1}^2|x^2-x|\,dx\)，说明各自选择换元、分部或区间拆分的依据。

（11）设\(f\)在\([0,1]\)上连续可导。证明\(\int_0^1f(x)\,dx=f(1)-\int_0^1xf'(x)\,dx\)。若\(f(0)=0\)、\(0\le f'\le1\)，估计\(\int_0^1 f\)。

（12）用四等分的中点与梯形公式估计\(\int_0^1dx/(1+x)\)。利用凸性判断上下界，再由\(\ln2\)核对。若只知道函数的单调性，能够得到哪些估计？

（13）细杆位于\([0,1]\)，线密度为\(\rho(x)=1+cx\)，其中\(c\ge0\)。求总质量和质心坐标，讨论质心怎样随\(c\)变化，并判断能否位于\(x=3/4\)。

（14）设\(f\)在\([0,1]\)上连续。证明存在\(\xi\in[0,1]\)，使\(\int_0^1xf(x)\,dx=f(\xi)/2\)。再说明把权函数\(x\)换成可能变号的函数时，结论为何未必成立。

（15）设\(f\)在\([0,1]\)上有界，且只有有限个间断点。它的累积函数是否一定连续？是否一定处处可导？分别说明依据并给出适当例子。

（16）分别构造：有无穷多个间断点但可积的函数；有界却不可积的函数；非零但绝对值积分为零的可积函数。比较这些例子与连续函数类的性质。

（17）（2025年北京卷第20题）已知函数\(f\)的定义域为\((-1,+\infty)\)，\(f(0)=0\)，且\(f'(x)=\ln(1+x)/(1+x)\)。设\(l_1\)是曲线在点\(A(a,f(a))\)、\(a\ne0\)处的切线。

\quad \exercisesubitem{1}求\(f'(x)\)的最大值；\exercisesubitem{2}当\(-1<a<0\)时，证明除切点外，曲线在\(l_1\)上方；\exercisesubitem{3}设过\(A\)且垂直于\(l_1\)的直线为\(l_2\)，\(l_1,l_2\)与横轴交点的横坐标分别为\(x_1,x_2\)。若\(a>0\)，求\((2a-x_1-x_2)/(x_2-x_1)\)的取值范围。

（18）（2022年新高考I卷第22题）已知函数\(f(x)=e^x-ax\)和\(g(x)=ax-\ln x\)有相同的最小值。

\quad \exercisesubitem{1}求\(a\)；\exercisesubitem{2}证明存在直线\(y=b\)，其与两条曲线\(y=f(x)\)、\(y=g(x)\)共有三个不同的交点，并且从左到右的三个交点的横坐标成等差数列。

（19）（2025年新课标II卷第18题）已知\(f(x)=\ln(1+x)-x+x^2/2-kx^3\)，其中\(0<k<1/3\)。

\quad \exercisesubitem{1}证明\(f\)在\((0,+\infty)\)内存在唯一的极值点和唯一的零点；\exercisesubitem{2}记它们分别为\(x_1,x_2\)，令\(g(t)=f(x_1+t)-f(x_1-t)\)。证明\(g\)在\((0,x_1)\)上单调递减，并比较\(2x_1\)与\(x_2\)的大小，证明你的结论。

（20）（2025年新课标I卷第19题）\exercisesubitem{1}求\(f(x)=5\cos x-\cos5x\)在\([0,\pi/4]\)上的最大值；\exercisesubitem{2}给定\(\theta\in(0,\pi)\)和\(a\in\mathbb R\)，证明存在\(y\in[a-\theta,a+\theta]\)，使\(\cos y\le\cos\theta\)；\exercisesubitem{3}设\(b\in\mathbb R\)，若存在\(\varphi\in\mathbb R\)，使\(5\cos x-\cos(5x+\varphi)\le b\)对所有实数\(x\)恒成立，求\(b\)的最小值。

\end{exercise}

\FloatBarrier
\newpage
\begin{exercise}[选做]
\mbox{}\par\nobreak

（1）设\(f(a)=a\)、\(f'(a)=\lambda\)，其中\(|\lambda|<1\)。由导数定义证明，存在\(q<1\)及\(a\)的一个邻域，使\(|f(x)-a|\le q|x-a|\)。据此研究从邻域内出发的迭代\(x_{n+1}=f(x_n)\)，给出收敛与误差估计，接回第一章的迭代数列。

（2）对\(f(x)=\ln(1+x)\)，用二阶Taylor公式估计\(\ln1.1\)，并给出误差上界。再用\(\int_1^{1.1}dt/t\)的中点公式估计同一个数，比较两种误差控制所需的信息。

（3）设\(f_n(x)=x^n\)，\(0\le x\le1\)。求逐点极限，证明不一致收敛，并计算\(\lim_n\int_0^1 f_n\)。这个例子能否说明“一致收敛是积分极限关系成立的必要条件”？

（4）设\(f_n\)是\([0,1]\)上的连续非负函数，\(\int_0^1f_n\to0\)。是否必有\(\|f_n\|_\infty\to0\)？尝试构造高度不变、宽度缩小的三角形图像，并解释面积误差与最大误差的区别。

（5）设\(F\in C^1[0,1]\)。证明\(F(x)=F(0)+\int_0^xF'(t)\,dt\)，并估计\(\|F\|_\infty\)。对\(F,G\)写出对应差值的估计，说明同时控制初值与导函数为什么能控制整个函数。

（6）用极坐标面积公式计算\(r=1+\cos\theta\)围成的面积，再把它写成平面参数曲线，用参数面积公式核对。说明两种局部贡献的表达为何不同。

（7）讨论\(\int_1^{+\infty}dx/[x(1+x^p)]\)的敛散性，其中\(p\in\mathbb R\)。收敛时给出尾部估计，再用积分判别法研究相应级数。

（8）设连续函数\(f_n,f\)定义在\([0,+\infty)\)，在每个有限区间上\(f_n\)一致收敛于\(f\)，且\(|f_n(x)|\le e^{-x}\)。证明\(f\)的反常积分绝对收敛，且\(\int_0^{+\infty}f_n\to\int_0^{+\infty}f\)。先固定截断位置控制有限范围，再统一估计尾部。

（9）设\(f\)在\([a,b]\)上有界、只有有限个间断点。用Darboux判据证明可积性，并说明“间断点附近缩小总长度”与“连续部分统一控制振幅”分别解决了什么问题。

（10）设\(f\)二阶连续可导，\(m\le f''\le M\)。用积分余项证明图像与任意切线之间的二次误差界，再由它推导中点数值积分在各小段上的误差界。
\end{exercise}
